% !PRINTED 651
\weberpage{651}

%<<CONT
in. One obtains in this way, by a simple calculation, the equation
\weq{11}{x^4 - 2\,(B^2c + b)\,x^2 - 4\,B\,a\,c\,x + B^4c^2 - 2\,B^2bc + c = 0.}
The representation of the $k$ by the formulae (9) is, however, not yet completely satisfactory, because among the rational quantities $a$, $b$, $c$ occurring in them the relation (4) still holds, and we seek a representation by independent quantities.

One special case we must first dispose of, namely $b = 0$, which has $a = i$ as a consequence. Then the formulae (9) give
\[
\begin{aligned}
k_0 &= C + B\sqrt{c} + \frac{1+i}{\sqrt{2}}\sqrt[4]{c}\\
k_1 &= C - B\sqrt{c} + \frac{1-i}{\sqrt{2}}\sqrt[4]{c}\\
k_2 &= C + B\sqrt{c} - \frac{1+i}{\sqrt{2}}\sqrt[4]{c}\\
k_3 &= C - B\sqrt{c} - \frac{1-i}{\sqrt{2}}\sqrt[4]{c},
\end{aligned}
\]
and the biquadratic equation (11) becomes
\[
x^4 - 2\,B^2c\,x^2 - 4\,i\,B\,c\,x + B^4c^2 + c = 0.
\]
But this case belongs here only when $i$ is contained in the field $\mathfrak{R}$, and hence never for real fields.

If, however, $b$ does not vanish, we introduce a new rational quantity $h$ by setting
\[
b = h\,(1 + a^2), \quad c = h^2\,(1 + a^2)
\]
whereby the relation (4) is then identically satisfied, and the formulae (9), if one replaces $Bh$ by $B$, give
\weq{12}{
\begin{aligned}
k_0 &= C + B\sqrt{1+a^2} + \sqrt{h\,\bigl(1+a^2+a\sqrt{1+a^2}\bigr)}\\
k_1 &= C - B\sqrt{1+a^2} + \sqrt{h\,\bigl(1+a^2-a\sqrt{1+a^2}\bigr)}\\
k_2 &= C + B\sqrt{1+a^2} - \sqrt{h\,\bigl(1+a^2+a\sqrt{1+a^2}\bigr)}\\
k_3 &= C - B\sqrt{1+a^2} - \sqrt{h\,\bigl(1+a^2-a\sqrt{1+a^2}\bigr)}
\end{aligned}
}
This expression for $k_0$ has moreover the advantage that, however one may determine the signs of the square roots occurring in it, it represents only four distinct values.
